\section{Related Work}
\label{sec:related}

Several systems combine static analysis with LLMs; our design was positioned against them before implementation.
\begin{itemize}
\item \textbf{LLMDFA}~\cite{llmdfa} extracts sources and sinks, summarizes functions with an LLM and validates path feasibility with external tools. It is a general LLM data-flow analysis, whereas we keep the symbolic engine as the only source of candidates.
\item \textbf{IRIS}~\cite{iris} uses LLMs to infer taint specifications for CodeQL and to filter its alerts in context. We instead freeze the specification and study triage alone.
\item \textbf{ZeroFalse}~\cite{zerofalse} is the closest: it enriches static alerts with flow context and CWE-specific prompts for structured LLM adjudication. Our difference is the deterministic gate that checks the model's cited evidence against extracted facts, the separate deterministic-only arm, and the separate accounting of semantic judgments.
\item \textbf{LLM4PFA}~\cite{llm4pfa} performs iterative interprocedural path-feasibility reasoning with LLM-generated SMT queries; we deliberately stay intraprocedural and solver-free.
\item \textbf{Sifting the Noise}~\cite{sifting2026} compares LLM agents for false-positive filtering and shows that they also discard genuine vulnerabilities. This motivates our retention-first metrics.
\item \textbf{PySeType}~\cite{pysetype} integrates LLM-assisted semantic type checking with Python vulnerability analysis.
\item \textbf{Memoir}~\cite{memoir} learns and verifies reusable false-positive patterns.
\end{itemize}
We use neither semantic types nor memory or retrieval. Our evaluation subjects come from the OWASP Benchmark for Python~\cite{owaspbench} and PyVul~\cite{pyvul}.
